\documentclass[aspectratio=169]{beamer}
\usetheme{metropolis}
\usecolortheme{default}

% =========================================
%               Packages
% =========================================
\usepackage[T1]{fontenc}
\usepackage[utf8]{inputenc}
\usepackage{graphicx}
\usepackage{amsmath, amssymb}
\usepackage{xcolor}
\usepackage{booktabs}
\usepackage{hyperref}
\usepackage{tikz}
\usepackage{tcolorbox}
\tcbuselibrary{skins, breakable}
\usepackage{microtype}
\SetTracking{encoding=*}{40}

% =========================================
%           Unicode Compatibility
% =========================================
\usepackage{newunicodechar}
\newunicodechar{→}{$\rightarrow$}
\newunicodechar{—}{\textemdash}
\newunicodechar{≠}{$\neq$}
\newunicodechar{≪}{$\ll$}

% =========================================
%             Custom Font Setup
% =========================================
\usepackage{FiraSans}
\renewcommand*\familydefault{\sfdefault}

% Slide Frame Title
\setbeamercolor{frametitle}{bg=white, fg=black}
\setbeamerfont{frametitle}{series=\bfseries, size=\Large}

% Title
\setbeamerfont{title}{size=\Large, series=\bfseries}

% =========================================
%        Highlight Color
% =========================================
\definecolor{SejongBlue}{RGB}{0,85,170}

% =========================================
%             Custom Boxes
% =========================================
% [MODIFIED] Applied the mybox style from the second text file.
\newtcolorbox{mybox}[2][]{
  colback=white,
  colframe=SejongBlue,
  fonttitle=\bfseries,
  title=#2,
  enhanced,
  drop shadow,
  arc=2mm,
  #1
}

\renewcommand{\textbf}[1]{\textcolor{SejongBlue}{\sf\bfseries #1}}

% =========================================
%                 Watermark
% =========================================
\usebackgroundtemplate{
  \tikz[overlay, remember picture]
  \node[opacity=0.06, at=(current page.center)] {
    \includegraphics[width=0.45\paperwidth]{1-1.교표.jpg}
  };
}

% =========================================
%             Title Information
% =========================================
\title{}
\author{}
\institute{}
\date{}

% =========================================
%                 Document
% =========================================
\begin{document}

% =========================================
%         CUSTOM TITLE PAGE FRAME
% =========================================
\begin{frame}[plain]

% --- Title ---
\vspace{0.7em} % [MODIFIED] 제목 위 여백 추가
{\LARGE \textcolor{SejongBlue}{\textbf{On Cross-Entropy under Class Imbalance:}}}\\[0.55em]
{\Large Stability and Weighting Perspectives}

% --- Conference Info ABOVE the line ---
\vspace{1.2em} % [MODIFIED] 제목과 학회 정보 간격 조정

\begin{center}
  {\small \textbf{International Workshop on Advanced Neural Networks}}\\[0.35em]
  {\small From Uncertainty-Aware Models to Hypercomplex Learning}\\[0.6em]
  {\footnotesize April 24\textsuperscript{th}, 2026}\\[0.25em]
\end{center}

% --- Horizontal Line ---
\vspace{1.0em}
{\color{SejongBlue}\hrule height 0.7pt}

% --- Authors CENTER aligned ---
\vspace{0.8em}
\begin{center}
{\Large
\textbf{Hyunseok Lee}$^{1}$ \quad
\textbf{Seung Ho Go}$^{2}$ \quad
\textbf{Jin Hee Yoon}$^{3}$
}\\[0.8em]

{\small % [MODIFIED] normlasize -> small
$^{1}$Department of Data Science \hspace{1.0em}
$^{2}$Department of Software \hspace{1.0em}
$^{3}$Department of AI \& IT\\[0.35em]
Sejong University, Korea
}
\end{center}

\end{frame}

% -------------------- Contents (All Sections) --------------------
\begin{frame}{Contents}
  \tableofcontents
\end{frame}

% Highlight current section in TOC (blue)
\setbeamercolor{section in toc current}{fg=SejongBlue}
\setbeamerfont{section in toc current}{series=\bfseries}


% -------------------- Section-wise TOC (Current Section Highlight) --------------------
\AtBeginSection[]{
  \begin{frame}{Contents}
    \tableofcontents[currentsection]
  \end{frame}
}

% =========================================
%   Section: Introduction
% =========================================
\section{Introduction}

% =========================================
%   Slide: Definition + Examples
% =========================================
\begin{frame}{Introduction: What is Class Imbalance?}

\footnotesize
\vspace{-0.8em}

\begin{columns}[T,totalwidth=\textwidth]

% ===== Left Box: Definition =====
\column{0.48\textwidth}
\begin{mybox}{Definition} % [MODIFIED] Converted from tcolorbox to mybox
\begin{itemize}
  \item Class imbalance occurs when some classes are \textbf{much rarer}.
  \item Rare classes receive fewer updates → often ignored by training.
\end{itemize}
\end{mybox}

% ===== Right Box: Examples =====
\column{0.48\textwidth}
\begin{mybox}{Examples} % [MODIFIED] Converted from tcolorbox to mybox
\begin{itemize}
  \item Fraud detection : 1 per 10,000 transactions
  \item Rare diseases : <1\%
  \item Defect detection : occurs in tens of thousands
\end{itemize}
\end{mybox}

\end{columns}

\vspace{1.2em}
\centering
{\Large \textit{``Rare classes define the real-world risk''}}

\end{frame}

% =========================================
%   Section: Existing Loss Functions
% =========================================
\section{Existing Loss Functions}

% --- Standard CE ---
\begin{frame}{Existing Loss Functions: Standard CE}

\begin{columns}[T, totalwidth=\textwidth]

% ---- Left: Text ----
\column{0.55\textwidth}
{\large \textbf{Standard Cross Entropy}}

{\large
\[
\mathcal{L}_{CE}(y,p) = - \sum_{c} y_c \log p_c
\]
}

{\large
\begin{itemize}
  \item No class weighting → majority dominates training.
  \item Minority gradients almost vanish.
\end{itemize}
}

{\color{gray}\footnotesize Issue: No minority learning signal.}

% ---- Right: Image ----
\column{0.45\textwidth}
\begin{center}
\includegraphics[width=1.20\linewidth]{standard_ce_distribution.png}
\end{center}

\end{columns}
\end{frame}

% --- Weighted CE ---
\begin{frame}{Existing Loss Functions: Weighted CE}

\begin{columns}[T, totalwidth=\textwidth]

% ---- Left: Text ----
\column{0.55\textwidth}
{\large \textbf{Weighted Cross Entropy}}

{\large
\[
w_c = \frac{1}{n_c}
\qquad\text{or}\qquad
w_c = \frac{1}{\sqrt{n_c}}
\]
\[
\mathcal{L}_{WCE} = - w_c \log p_c
\]
}

{\large
\begin{itemize}
  \item Boosts rare-class gradients via inverse-frequency.
  \item Overshoots under extreme imbalance.
\end{itemize}
}

{\color{gray}\footnotesize Issue: Weight explosion.}

% ---- Right: Image ----
\column{0.45\textwidth}
\begin{center}
\includegraphics[width=1.20\linewidth]{weighted_ce_explosion.png}
\end{center}

\end{columns}
\end{frame}

% --- Focal Loss ---
\begin{frame}{Existing Loss Functions: Focal Loss}

\begin{columns}[T,totalwidth=\textwidth]

% ---- Left: Text ----
\column{0.55\textwidth}
{\large \textbf{Focal Loss}}

{\large
\[
\mathcal{L}_{focal} = - \alpha_t (1 - p_t)^\gamma \log(p_t)
\]

}

{\large
\begin{itemize}
  \item Down-weights easy samples.
  \item Improves minority recall with hard-sample focus.
  \item Includes $\alpha$ for class balancing, $\gamma$ controls modulating strength
\end{itemize}
}

{\color{gray}\footnotesize Issue: Extremely sensitive to $\gamma, \alpha$.}

% ---- Right: Image ----
\column{0.45\textwidth}
\begin{center}
\includegraphics[width=1.20\linewidth]{focal_loss_curves.png}
\end{center}

\end{columns}
\end{frame}

% --- CBCE ---
\begin{frame}{Existing Loss Functions: Class-Balanced Loss}

\begin{columns}[T,totalwidth=\textwidth]

% ---- Left: Text ----
\column{0.55\textwidth}
{\large \textbf{Class-Balanced Loss}}

{\large
\[
E(n_c)=\frac{1-\beta^{n_c}}{1-\beta}, \qquad
w_c=\frac{1}{E(n_c)}
\]
}

{\large
\begin{itemize}
  \item Uses “effective number of samples”.
  \item More stable than WCE.
\end{itemize}
}

{\color{gray}\footnotesize Issue: Still unstable for ultra-rare classes.}

% ---- Right: Image ----
\column{0.45\textwidth}
\begin{center}
\includegraphics[width=1.20\linewidth]{cbce_effective_number.png}
\end{center}

\end{columns}
\end{frame}

% =========================================
%             Limitations
% =========================================
\begin{frame}{Existing Loss Functions: Limitations}

\textbf{Common Problems}

\begin{itemize}
  \item Hand-crafted weights → inconsistent behavior
  \item No smooth scaling for extremely small classes
  \item Uncontrolled gradients → unstable convergence
\end{itemize}

\vspace{0.6em}
\centering
\textit{→ Motivates ALW-CE and ES-ACE}

\end{frame}

% =========================================
%   Section: Proposed Methods
% =========================================
\section{Proposed Methods}

% -------------------- ALW-CE --------------------
\begin{frame}{Proposed: Adaptive Log-Weighted CE (ALW-CE)}

\begin{columns}[T,totalwidth=\textwidth]

% ----- Left Text -----
\column{0.55\textwidth}

\textbf{Weight definition}
\[
w_c = \frac{1}{\log(1 + n_c)}
\]

\textbf{Interpretation}
\begin{itemize}
  \item \( n_c \): sample count of class \( c \)
  \item Log scaling prevents extreme weights.
\end{itemize}

\textbf{Limitation}
\begin{itemize}
  \item When \(n_c\) is extremely small, \(\log(1+n_c)\) is still tiny.
  \item Weight can still grow excessively.
\end{itemize}

% ----- Right Graph -----
\column{0.45\textwidth}
\begin{center}
\includegraphics[width=1.05\linewidth]{alw_ce_curve_smooth.png}
\end{center}

\end{columns}

\end{frame}

% -------------------- PLWCE [ADDED] --------------------
% [MODIFIED] Split into two frames for better readability.

\begin{frame}{Related Extension: PLWCE}

\vspace{-0.3em}

\begin{columns}[T,totalwidth=\textwidth]

% ----- Left Text -----
\column{0.60\textwidth}

\textbf{Power-Scaled Log-Weighted CE (PLWCE)}
\[
w_c = \frac{1}{(\log(1+n_c))^{\alpha}}
\]

\textbf{Interpretation}
\begin{itemize}
  \item Extends ALW-CE with a tunable exponent $\alpha$.
  \item Provides dataset-specific control of weighting strength.
\end{itemize}

\textbf{Effect of $\alpha$}
\begin{itemize}
  \item $\alpha > 1$: stronger emphasis on minority classes
  \item $\alpha < 1$: more conservative weighting
\end{itemize}

\vspace{0.3em}
{\footnotesize \textcolor{gray}{Explored as a related extension in our joint research.}}

% ----- Right Summary Box -----
\column{0.38\textwidth}
\begin{mybox}[before skip=0pt, after skip=0pt]{Key Message}
\small
PLWCE introduces a tunable exponent $\alpha$ on top of the log-weighted formulation.

\vspace{0.4em}
This makes the weighting strategy more flexible across different imbalance settings.
\end{mybox}

\end{columns}

\end{frame}

% -------------------- PLWCE vs ES-ACE [ADDED] --------------------
% [MODIFIED] Comparison moved to a separate frame.

\begin{frame}{Two Extensions of Log-Weighted CE}

\vspace{-0.5em}

\begin{columns}[T,totalwidth=\textwidth]

% ----- Left: PLWCE -----
\column{0.48\textwidth}
\begin{mybox}[before skip=0pt, after skip=0pt, top=2pt, bottom=2pt]{PLWCE}
\[
w_c = \frac{1}{(\log(1+n_c))^{\alpha}}
\]

\textbf{Goal}
\begin{itemize}
  \item Control weighting intensity
\end{itemize}

\textbf{Key parameter}
\begin{itemize}
  \item $\alpha$
\end{itemize}

\textbf{Main effect}
\begin{itemize}
  \item Adjustable focus on minority classes
\end{itemize}
\end{mybox}

% ----- Right: ES-ACE -----
\column{0.48\textwidth}
\begin{mybox}[before skip=0pt, after skip=0pt, top=2pt, bottom=2pt]{ES-ACE}
\[
w_c = \frac{1}{\log(1+n_c) + \varepsilon}
\]

\textbf{Goal}
\begin{itemize}
  \item Improve stability for ultra-rare classes
\end{itemize}

\textbf{Key parameter}
\begin{itemize}
  \item $\varepsilon$
\end{itemize}

\textbf{Main effect}
\begin{itemize}
  \item Prevent excessive weights and stabilize gradients
\end{itemize}
\end{mybox}

\end{columns}

\vspace{0.25em}
\centering
{\small \textit{PLWCE focuses on controllable emphasis, whereas ES-ACE focuses on stable smoothing.}}

\end{frame}

% -------------------- ES-ACE --------------------
\begin{frame}{Proposed: $\varepsilon$-Smoothed Adaptive Log-Weighted CE (ES-ACE)}

\begin{columns}[T,totalwidth=\textwidth]

% ----- Left Text -----
\column{0.55\textwidth}
\vspace{-0.5em}
% [ADDED] Transition sentence from PLWCE to ES-ACE.
{\small Alongside the $\alpha$-based PLWCE extension, we further explored an $\varepsilon$-smoothed variant for improved stability.}

\textbf{Weight definition}
\[
w_c = \frac{1}{\log(1+n_c) + \varepsilon}
\]

\textbf{Role of $\varepsilon$}
\begin{itemize}
  \item Prevents weight explosion for ultra-rare classes.
  \item Smooths the weight curve → more stable gradients.
  \item Reduces overfitting from extremely large rare-class weights.
\end{itemize}

% ----- Right Graph -----
\column{0.45\textwidth}
\begin{center}
\includegraphics[width=1.05\linewidth]{es_ace_curve_smooth.png}
\end{center}

\end{columns}

\end{frame}



% -------------------- Effect of epsilon range --------------------
\begin{frame}{Effect of $\varepsilon$ Range}
\centering

\small
\scalebox{0.90}{%
\begin{tabular}{ccc}
\toprule
\textbf{$\varepsilon$ Range} & \textbf{Behavior} & \textbf{Observed Performance} \\
\midrule
$\varepsilon \approx 0$
& Behaves almost like ALW-CE
& \textcolor{gray}{Stabilized but still limited for ultra-rare classes} \\
\midrule
$0.1 \sim 0.5$
& Mild smoothing, controlled weights
& \textcolor{blue}{Best overall stability–performance trade-off} \\
\midrule
$0.5 \sim 1.0$
& Stronger smoothing
& \textcolor{blue}{Stable, beneficial on noisy datasets} \\
\midrule
$> 1.0$
& Excessive smoothing
& \textcolor{gray}{No significant performance gain observed} \\
\bottomrule
\end{tabular}%
}
\end{frame}


% =========================================
%   Section: Experimental Setup
% =========================================
\section{Experimental Setup}

% -------------------- Experimental Setup --------------------
\begin{frame}{Experimental Setup}

\textbf{Datasets \& Model Setup}
\begin{itemize}
  \item 11 benchmark datasets (finance, biology, fault detection)
  \item Model: XGBoost (gblinear)
  \item Losses: CE, WCE, CB, Focal, ALW-CE, ES-ACE
\end{itemize}

\end{frame}

% -------------------- Experimental Setup (2) --------------------
\begin{frame}{Experimental Setup (2/2)}

\textbf{Hyperparameter Optimization}
\begin{itemize}
  \item Optuna: 30 trials × 3-fold CV
  \item Target metric: Macro F1-score
\end{itemize}

\textbf{Evaluation Metrics}
\begin{itemize}
  \item Macro F1 
  \item PR-AUC 
\end{itemize}

\end{frame}

% =========================================
%   Section: Experimental Results
% =========================================
\section{Experimental Results}

% =========================================
%   Full Experimental Results (I)
% =========================================
\begin{frame}{Full Experimental Results (I)}
\scriptsize
\renewcommand{\arraystretch}{0.85}

\begin{columns}[T, totalwidth=\textwidth]

% ============================
% Column 1
% ============================
\column{0.30\textwidth}

% ---- Credit Card Fraud ----
\textbf{Credit Card Fraud}
\begin{tabular}{lrr}
\toprule
Loss & F1 & PR \\
\midrule
CE & 0.7287 & 0.7073 \\
WCE & 0.7917 & 0.6984 \\
WCE(sq) & \textcolor{SejongBlue}{\textbf{0.7972}} & 0.6996 \\
Focal & 0.0037 & 0.6542 \\
CB & 0.7883 & 0.6999 \\
\textcolor{SejongBlue}{ALW-CE} & 0.7447 & \textcolor{red}{\textbf{0.7155}} \\
\textcolor{SejongBlue}{ES-ACE} & 0.6988 & 0.7092 \\
\bottomrule
\end{tabular}

\vspace{0.8em}

% ---- German Credit ----
\textbf{German Credit}
\begin{tabular}{lrr}
\toprule
Loss & F1 & PR \\
\midrule
CE & 0.5854 & 0.6391 \\
WCE & \textcolor{SejongBlue}{\textbf{0.6301}} & 0.6321 \\
WCE(sq) & 0.6146 & 0.6446 \\
Focal & 0.5138 & 0.6215 \\
CB & 0.6070 & \textcolor{red}{\textbf{0.6450}} \\
\textcolor{SejongBlue}{ALW-CE} & 0.5909 & 0.6398 \\
\textcolor{SejongBlue}{ES-ACE} & 0.5802 & 0.6397 \\
\bottomrule
\end{tabular}

% ============================
% Column 2
% ============================
\column{0.30\textwidth}

% ---- Aps Failure ----
\textbf{Aps Failure}
\begin{tabular}{lrr}
\toprule
Loss & F1 & PR \\
\midrule
CE & 0.7595 & 0.8075 \\
WCE & 0.6447 & 0.8107 \\
WCE(sq) & 0.7581 & 0.8242 \\
Focal & 0.0170 & 0.0210 \\
CB & \textcolor{SejongBlue}{\textbf{0.7688}} & \textcolor{red}{\textbf{0.8253}} \\
\textcolor{SejongBlue}{ALW-CE} & 0.7662 & 0.8246 \\
\textcolor{SejongBlue}{ES-ACE} & 0.7518 & 0.8018 \\
\bottomrule
\end{tabular}

\vspace{0.8em}

% ---- Bank Marketing ----
\textbf{Bank Marketing}
\begin{tabular}{lrr}
\toprule
Loss & F1 & PR \\
\midrule
CE & 0.4533 & 0.5482 \\
WCE & 0.5514 & 0.5413 \\
WCE(sq) & \textcolor{SejongBlue}{\textbf{0.5735}} & 0.5455 \\
Focal & 0.2652 & 0.4039 \\
CB & 0.4752 & 0.5034 \\
\textcolor{SejongBlue}{ALW-CE} & 0.4889 & 0.5476 \\
\textcolor{SejongBlue}{ES-ACE} & 0.4565 & \textcolor{red}{\textbf{0.5482}} \\
\bottomrule
\end{tabular}

% ============================
% Column 3
% ============================
\column{0.30\textwidth}

% ---- Secom ----
\textbf{Secom}
\begin{tabular}{lrr}
\toprule
Loss & F1 & PR \\
\midrule
CE & 0.1667 & 0.1396 \\
WCE & \textcolor{SejongBlue}{\textbf{0.2759}} & \textcolor{red}{\textbf{0.1748}} \\
WCE(sq) & 0.1695 & 0.1370 \\
Focal & 0.1469 & 0.1632 \\
CB & 0.2264 & 0.1746 \\
\textcolor{SejongBlue}{ALW-CE} & 0.2000 & 0.1689 \\
\textcolor{SejongBlue}{ES-ACE} & 0.1667 & 0.1403 \\
\bottomrule
\end{tabular}

\vspace{0.8em}

% ---- Credit Card Default ----
\textbf{Credit Card Default}
\begin{tabular}{lrr}
\toprule
Loss & F1 & PR \\
\midrule
CE & 0.3568 & 0.4982 \\
WCE & \textcolor{SejongBlue}{\textbf{0.5023}} & 0.4663 \\
WCE(sq) & 0.4938 & 0.4839 \\
Focal & 0.3859 & 0.4625 \\
CB & 0.4873 & 0.4714 \\
\textcolor{SejongBlue}{ALW-CE} & 0.4913 & 0.4726 \\
\textcolor{SejongBlue}{ES-ACE} & 0.3568 & \textcolor{red}{\textbf{0.4982}} \\
\bottomrule
\end{tabular}

\end{columns}
\end{frame}

% =========================================
%   Full Experimental Results (II)
% =========================================
\begin{frame}{Full Experimental Results (II)}
\scriptsize
\renewcommand{\arraystretch}{0.85}

\begin{columns}[T, totalwidth=\textwidth]

% ============================
% Column 1
% ============================
\column{0.30\textwidth}

% ---- Telco Churn ----
\textbf{Telco Churn}
\begin{tabular}{lrr}
\toprule
Loss & F1 & PR \\
\midrule
CE & 0.5942 & 0.6399 \\
WCE & 0.6225 & 0.6406 \\
WCE(sq) & 0.6252 & 0.6287 \\
Focal & 0.5402 & 0.6184 \\
CB & \textcolor{SejongBlue}{\textbf{0.6295}} & 0.6275 \\
\textcolor{SejongBlue}{ALW-CE} & 0.6286 & 0.6285 \\
\textcolor{SejongBlue}{ES-ACE} & 0.6066 & \textcolor{red}{\textbf{0.6422}} \\
\bottomrule
\end{tabular}

\vspace{0.8em}

% ---- Yeast ----
\textbf{Yeast (Multi)}
\begin{tabular}{lrr}
\toprule
Loss & F1 & PR \\
\midrule
CE & 0.5843 & 0.5964 \\
WCE & 0.4591 & 0.5242 \\
WCE(sq) & 0.5212 & 0.5934 \\
Focal & 0.4786 & 0.5132 \\
CB & 0.5530 & 0.5975 \\
\textcolor{SejongBlue}{ALW-CE} & 0.5445 & \textcolor{red}{\textbf{0.5980}} \\
\textcolor{SejongBlue}{ES-ACE} & \textcolor{SejongBlue}{\textbf{0.5714}} & 0.5969 \\
\bottomrule
\end{tabular}

% ============================
% Column 2
% ============================
\column{0.30\textwidth}

% ---- Page Blocks ----
\textbf{Page Blocks (Multi)}
\begin{tabular}{lrr}
\toprule
Loss & F1 & PR \\
\midrule
CE & 0.8292 & 0.8855 \\
WCE & 0.7226 & 0.8587 \\
WCE(sq) & 0.8249 & 0.8792 \\
Focal & 0.7735 & 0.8477 \\
CB & \textcolor{SejongBlue}{\textbf{0.8304}} & 0.8782 \\
\textcolor{SejongBlue}{ALW-CE} & 0.8153 & 0.8916 \\
\textcolor{SejongBlue}{ES-ACE} & 0.8153 & \textcolor{red}{\textbf{0.8922}} \\
\bottomrule
\end{tabular}

\vspace{0.8em}

% ---- Glass ----
\textbf{Glass (Multi)}
\begin{tabular}{lrr}
\toprule
Loss & F1 & PR \\
\midrule
CE & 0.6380 & 0.7269 \\
WCE & 0.6734 & 0.7417 \\
WCE(sq) & 0.6972 & 0.7309 \\
Focal & 0.6568 & 0.7307 \\
CB & 0.7095 & 0.7383 \\
\textcolor{SejongBlue}{ALW-CE} & \textcolor{SejongBlue}{\textbf{0.7361}} & \textcolor{red}{\textbf{0.7391}} \\
\textcolor{SejongBlue}{ES-ACE} & 0.6996 & 0.7384 \\
\bottomrule
\end{tabular}

% ============================
% Column 3
% ============================
\column{0.30\textwidth}

% ---- Steel Faults ----
\textbf{Steel Faults (Multi)}
\begin{tabular}{lrr}
\toprule
Loss & F1 & PR \\
\midrule
CE & 0.6902 & 0.7270 \\
WCE & 0.6316 & 0.7193 \\
WCE(sq) & 0.6949 & 0.7249 \\
Focal & 0.6569 & 0.6880 \\
CB & 0.6997 & \textcolor{red}{\textbf{0.7322}} \\
\textcolor{SejongBlue}{ALW-CE} & \textcolor{SejongBlue}{\textbf{0.7075}} & 0.7256 \\
\textcolor{SejongBlue}{ES-ACE} & 0.6996 & 0.7262 \\
\bottomrule
\end{tabular}

\end{columns}
\end{frame}

% -------------------- PLWCE RESULT SUMMARY (1/2) --------------------
% [MODIFIED] 기존 1장짜리 결과 요약 슬라이드를 2장으로 분리

\begin{frame}{Related Experimental Finding: PLWCE}

\vspace{-0.4em} % [MODIFIED] 제목 아래 간격 축소

\begin{mybox}[before skip=0pt, after skip=0pt, top=2pt, bottom=2pt]{Why mention PLWCE in the results?} % [MODIFIED]
\small
Since PLWCE was explored as a related extension in our joint research,
we briefly summarize its empirical behavior here.
\end{mybox}

\vspace{2em} % [MODIFIED]

\begin{mybox}[before skip=0pt, after skip=0pt, top=2pt, bottom=2pt]{Key Observation} % [MODIFIED]
\small
PLWCE, the $\alpha$-based extension of log-weighted CE, showed
strong F1 improvements on several highly imbalanced datasets.

\vspace{0.4em}

In the previous joint study, PLWCE achieved higher F1 than LWCE
in 5 out of 11 datasets, showing that tunable minority emphasis
can be beneficial in some settings.
\end{mybox}

\end{frame}

% -------------------- PLWCE RESULT SUMMARY (2/2) --------------------
% [MODIFIED] Takeaway를 별도 슬라이드로 분리

\begin{frame}{Related Experimental Finding: PLWCE}

\vspace{-0.4em} % [MODIFIED]

\begin{mybox}[before skip=0pt, after skip=0pt, top=2pt, bottom=2pt]{Takeaway} % [MODIFIED]
\small
PLWCE and ES-ACE should be viewed as complementary extensions:
\begin{itemize}
  \item \textbf{PLWCE}: improves controllability via $\alpha$
  \item \textbf{ES-ACE}: improves stability via $\varepsilon$
\end{itemize}
\end{mybox}

\vspace{2em} % [MODIFIED]

\begin{mybox}[before skip=0pt, after skip=0pt, top=2pt, bottom=2pt]{Interpretation} % [MODIFIED]
\small
This means the two methods do not compete in exactly the same direction.

\vspace{0.4em}

PLWCE is useful when stronger and more flexible minority emphasis is needed,
whereas ES-ACE is useful when stable smoothing is more important.
\end{mybox}

\end{frame}

% =========================================
%   Borda Summary Table (F1 & PR-AUC)
% =========================================
\begin{frame}{Borda Summary: F1 \& PR-AUC Rankings}

\centering
\small
\renewcommand{\arraystretch}{1.15}

\begin{tabular}{lcccc}
\toprule
\textbf{Loss} &
\textbf{Avg Rank (F1)} & \textbf{Borda (F1)} &
\textbf{Avg Rank (PR)} & \textbf{Borda (PR)} \\
\midrule
es\_ace               & 4.36 & 40 & \textbf{2.73} & \textbf{58} \\
alw\_ce          & 3.50 & 49 & 3.00 & 55 \\
class\_balanced\_loss   & \textbf{2.64} & \textbf{59} & 3.55 & 49 \\
standard\_ce          & 4.32 & 40 & 3.55 & 49 \\
weighted\_ce\_inv     & 4.09 & 43 & 5.09 & 32 \\
weighted\_ce\_inv\_sqrt & 3.00 & 55 & 3.91 & 45 \\
focal\_loss           & 6.09 & 21 & 6.18 & 20 \\
\bottomrule
\end{tabular}

\vspace{0.7em}

{\footnotesize
\textcolor{gray}{\textit{Lower Avg Rank is better. Higher Borda score is better.}}
}

\end{frame}


% -------------------- F1 Results --------------------
\begin{frame}{Results: Macro F1-Score}

\textbf{Why ES-ACE F1 is lower?}
\begin{itemize}
  \item $\varepsilon$-smoothing reduces extreme class weights.
  \item Prevents overfitting but slightly lowers recall for moderate imbalance.
\end{itemize}

\begin{center}
  \includegraphics[width=0.55\linewidth]{borda_summary_F1-Score_e3.png}
\end{center}

\end{frame}

% -------------------- PR-AUC Results --------------------
\begin{frame}{Results: PR-AUC}

\textbf{Why ES-ACE excels in PR-AUC?}
\begin{itemize}
  \item PR-AUC reflects minority recall under extreme imbalance.
  \item $\varepsilon$-smoothing stabilizes gradients and avoids weight explosion.
  \item Enables reliable detection of ultra-rare classes.
\end{itemize}

\begin{center}
  \includegraphics[width=0.55\linewidth]{borda_summary_PR_AUC_e3.png}
\end{center}

\end{frame}

% =========================================
%   Section: Conclusion
% =========================================
\section{Conclusion}

% =========================================
%   Discussion (1/2)
% =========================================
\begin{frame}{Conclusion}
\small

\begin{mybox}{Discussion} % [MODIFIED] Converted from tcolorbox to mybox

\begin{enumerate}

% 1 — Proposed ALW-CE
\item \textbf{Proposed ALW-CE:}
Adaptive Log-Weighted CE:
\[
w_c = \frac{1}{\log(1+n_c)}
\]
which mitigates inverse-frequency instability and provides
a smooth, more stable weighting baseline.

\vspace{0.4em}

% 2 — Proposed ES-ACE
\item \textbf{Proposed ES-ACE:}
Extended into $\varepsilon$-Smoothed ALW-CE:
\[
w_c = \frac{1}{\log(1+n_c) + \varepsilon}
\]
preventing weight explosion for extremely rare classes
and stabilizing gradients.

\end{enumerate}

\end{mybox}

\end{frame}

% =========================================
%   Discussion (2/2)
% =========================================
\begin{frame}{Conclusion}
\small
\vspace{-1em}
\begin{mybox}{Discussion} % [MODIFIED] Converted from tcolorbox to mybox

\begin{enumerate}
\setcounter{enumi}{2}

% 3 — Analyzed the Role of ε
\item \textbf{Analyzed the Role of $\varepsilon$:}
Optuna tuning showed that:
\[
0.1 \le \varepsilon \le 1.0
\]
provides the best stability–performance trade-off.

\vspace{0.4em}

% 4 — Validated Performance
\item \textbf{Validated Performance:}
Across 11 datasets, ES-ACE achieved consistently strong PR-AUC
and competitive F1, demonstrating robustness in real imbalance tasks.

\vspace{0.4em}

% [ADDED] Joint research note on PLWCE.
\item \textbf{Related Joint Extension:}
Our joint research also explored \textbf{PLWCE} as an $\alpha$-based extension,
while this presentation focuses on the $\varepsilon$-smoothed \textbf{ES-ACE} variant.

\end{enumerate}

\end{mybox}

\vspace{0.6em}
\centering
{\large \textit{“Minimal smoothing, substantial stability, practical performance.”}}

\end{frame}

% =========================================
%            References
% =========================================
\begin{frame}{References}
\small
\begin{enumerate}

\item T. H. Phan and K. Yamamoto,
\textit{“Resolving Class Imbalance in Object Detection with Weighted Cross Entropy Losses,”}
arXiv preprint arXiv:2006.01413, 2020.

\item T.-Y. Lin, P. Goyal, R. Girshick, K. He, and P. Dollár,
\textit{“Focal Loss for Dense Object Detection,”}
Proc. IEEE Int. Conf. Computer Vision (ICCV), pp. 2980–2988, 2017.

\item Y. Cui, M. Jia, T. Lin, Y. Song, and S. Belongie,
\textit{“Class-Balanced Loss Based on Effective Number of Samples,”}
Proc. IEEE/CVF Conf. Computer Vision and Pattern Recognition (CVPR),
pp. 9268–9277, 2019.

\item T. H. Phan and K. Yamamoto,
\textit{“Resolving Class Imbalance in Object Detection with Weighted Cross Entropy Losses,”}
arXiv preprint arXiv:2006.01413, 2020.

\end{enumerate}
\end{frame}


% =========================================
%            Final Slide
% =========================================
\begin{frame}[plain]
\centering

\vspace{1.5em}

{\Huge \textbf{Thank you for your attention!}}\\[1.2em]

{\Large Any Questions?}\\[2.2em]

% Contact Box
\begin{mybox}[width=0.65\textwidth, center, left=10pt, right=10pt, top=10pt, bottom=10pt]{Contact Information} % [MODIFIED] Converted from tcolorbox to mybox

\centering
{\Large \textbf{Hyunseok Lee}} \\[0.5em]
{\large \texttt{hslee901@gmail.com}} \\[0.3em]
{\large \texttt{010-2054-2689}} \\[0.3em]
{\large \texttt{GitHub: @HyunseokLee928}}

\end{mybox}

\vspace{1.2em}
{\small \textit{Sejong University · Department of Data Science}}

\end{frame}

\end{document}
m